% =====================================================================
% Vendor benchmark chapter template — STM32Cube Ecosystem Benchmark report
%
% The report compares ST (the reference) against ONE competing vendor per
% chapter (e.g. GD32, TI, NXP). Copy the relevant blocks into the target
% file, replace <...> placeholders, and delete unused blocks. Keep labels
% meaningful and unique. pdfLaTeX-safe.
%
% Chapter layout (fixed so every vendor reads the same way):
%   Introduction
%   -> Hardware platforms        (board-characteristics comparison table)
%   -> one \section per stack    (USB device stack, FreeRTOS, ...), each:
%        Scenario design -> Results (per scenario) -> Interpretation
%   -> Conclusion                (ST vs this vendor synthesis)
%
% Routing reminder:
%   new vendor          -> new chapter (this whole file, wired into main.tex)
%   existing vendor,
%     new stack         -> just a new \section block (the STACK block below)
%   existing stack,
%     new scenario      -> just a new \subsubsection in that stack's Results
% =====================================================================

% =========================== CHAPTER =================================
\chapter{Benchmarking ST against <Vendor, e.g. GD32>}
\label{chap:vs-<vendor-key>}

% --- Introduction: why this vendor, what the comparison settles ------
We compare the STMicroelectronics / STM32Cube ecosystem against
<vendor>'s equivalent offering on <which stacks, e.g. the USB device stack
and the FreeRTOS integration>. <Say why this vendor is a relevant competitor
(market position, ST-compatibility, price) and what question this comparison
should settle.>

% --- Hardware platforms: the two boards under test ------------------
\section{Hardware platforms}
\label{sec:<vendor-key>-boards}
Every scenario in this chapter runs on two comparable boards: the ST
reference board <ST board / part> and <vendor>'s <vendor board / part>.
Table~\ref{tab:<vendor-key>-boards} contrasts their characteristics.

\begin{table}[H]
  \centering
  \begin{tabular}{l l l}
    \toprule
    Characteristic   & ST (<ST part>)            & <Vendor> (<part>) \\
    \midrule
    Core / clock     & <Cortex-M? / freq\,MHz>   & <core / freq\,MHz> \\
    Flash / RAM      & <KB / KB>                 & <KB / KB> \\
    Key peripherals  & <USB FS/HS, ...>          & <...> \\
    Toolchain / IDE  & STM32CubeIDE / GCC        & <vendor IDE / GCC> \\
    Reference price  & <approx>                  & <approx> \\
    \bottomrule
  \end{tabular}
  \caption{Board characteristics: ST reference vs <vendor>.}
  \label{tab:<vendor-key>-boards}
\end{table}

% ====================== STACK BLOCK (repeat per stack) ===============
% Copy this whole \section for each benchmarked stack (USB, FreeRTOS, ...).
\section{<Stack, e.g. USB device stack>}
\label{sec:<vendor-key>-<stack-key>}

% --- Scenario design: which scenarios and why -----------------------
\subsection{Scenario design}
\label{sec:<vendor-key>-<stack-key>-scenarios}
To exercise the <stack> meaningfully we defined <N> scenarios:
<e.g. HID, CDC, and MSC>. <Justify the choice --- they cover the dominant
device classes / use cases.> Each scenario was implemented identically on
both platforms; Table~\ref{tab:<vendor-key>-<stack-key>-scenarios} states
what each one exercises and the metric it yields.

\begin{table}[H]
  \centering
  \begin{tabular}{l l l}
    \toprule
    Scenario & What it exercises          & Measured metric \\
    \midrule
    <HID>    & <low-rate interrupt xfers> & <latency / cycles> \\
    <CDC>    & <virtual COM throughput>   & <MB/s> \\
    <MSC>    & <mass-storage throughput>  & <MB/s> \\
    \bottomrule
  \end{tabular}
  \caption{<Stack> benchmark scenarios.}
  \label{tab:<vendor-key>-<stack-key>-scenarios}
\end{table}

% --- Results: one block per scenario --------------------------------
\subsection{Results}
\label{sec:<vendor-key>-<stack-key>-results}

\subsubsection{<Scenario 1, e.g. HID>}
<One line on any scenario-specific setup.> Table~\ref{tab:<vendor-key>-<stack-key>-s1}
reports <metric + units> for both platforms.

\begin{table}[H]
  \centering
  \begin{tabular}{l r r}
    \toprule
    Platform           & <Metric (unit)> & <Metric 2 (unit)> \\
    \midrule
    ST (<part>)        & <value>         & <value> \\
    <Vendor> (<part>)  & <value>         & <value> \\
    \bottomrule
  \end{tabular}
  \caption{<Scenario 1> results: ST vs <vendor>.}
  \label{tab:<vendor-key>-<stack-key>-s1}
\end{table}

% \subsubsection{<Scenario 2>}  ... repeat the results table ...
% \subsubsection{<Scenario 3>}  ... repeat the results table ...

% Optional: a grouped bar chart comparing all scenarios at once.
% \begin{figure}[H]
%   \centering
%   \includegraphics[width=0.8\textwidth]{Figures/<vendor-key>/<stack-key>-chart}
%   \caption{<Stack> throughput across scenarios: ST vs <vendor>.}
%   \label{fig:<vendor-key>-<stack-key>-chart}
% \end{figure}

% --- Interpretation: what the comparison reveals --------------------
\subsection{Interpretation}
<Across the scenarios, where does ST lead or trail <vendor>, and by how much?
Tie each difference back to the board characteristics
(Table~\ref{tab:<vendor-key>-boards}) or to the stack implementation. Name the
trade-off and the practical guidance for an engineer choosing between the two.
Cite any external spec used for context~\cite{<key>}.>

% ====================== END STACK BLOCK =============================
% (repeat the \section above for the next stack: FreeRTOS, ...)

% --- Chapter conclusion: ST vs this vendor --------------------------
\section*{Conclusion}
<Synthesize how the STM32Cube ecosystem stands against <vendor> across the
stacks tested, and the practical guidance for an engineer choosing between
them. Feed this takeaway into the optional cross-vendor synthesis chapter.>
